% ============================================================================
\chapter{PMM：物理内存管理}
\label{ch:pmm}
% Stage 5 — 对应 2026-03-11/12 changelog
% ============================================================================

\section{本章目标}

\begin{itemize}
  \item 解析 Multiboot2 Memory Map，获取物理内存布局
  \item 设计可插拔 PMM 后端接口（\texttt{pmm\_ops\_t}）
  \item 实现 bitmap 分配器和 buddy 分配器两个后端
  \item Shell 命令：\texttt{free}、\texttt{mmap}、\texttt{pmm}
\end{itemize}

\section{背景知识：物理内存布局}

BIOS 通过 Multiboot2 Memory Map Tag（type=6）报告物理内存区域：

\begin{lstlisting}[style=shellstyle]
Available RAM: 0x00000000 - 0x0009FC00 (640KiB)   # 低端常规内存
Reserved:      0x0009FC00 - 0x000FFFFF             # VGA/BIOS/ROM
Available RAM: 0x00100000 - 0x07FE0000 (126MiB)    # 高端可用内存
\end{lstlisting}

PMM 的任务：管理"Available RAM"中的物理页（每页 4KiB），提供 alloc/free 接口。

\begin{thinkbox}
为什么 PMM 以"页"（4KiB）为单位分配，而不是字节？提示：想想后续分页机制中
PTE 指向的最小单位是什么。
\end{thinkbox}

\section{可插拔后端架构}

\begin{lstlisting}[style=cstyle]
typedef struct pmm_ops {
    const char* name;
    void  (*init)(void);
    void* (*alloc_page)(void);      // 返回物理地址，NULL=OOM
    void  (*free_page)(void* page);
    unsigned (*total_pages)(void);
    unsigned (*free_pages)(void);
    unsigned (*managed_base)(void);
    // ...
} pmm_ops_t;
\end{lstlisting}

\texttt{pmm.c} 是薄 dispatch 层：所有调用通过 \texttt{g\_pmm\_ops} 转发。
切换后端只需在 \texttt{kmain} 中改一行 \texttt{pmm\_register\_backend()}。

\begin{designbox}
这套"ops 函数指针表 + dispatch 层"的模式参考了 Linux 内核的 \texttt{file\_operations}。
后续 kmalloc 和 VMA 也采用完全相同的架构，形成统一的可插拔设计风格。
\end{designbox}

\section{Bitmap 分配器}

最直观的方案：每个物理页对应 1 bit。

\begin{itemize}
  \item \texttt{alloc\_page}：线性扫描找第一个 0 bit → 置 1 → 返回物理地址
  \item \texttt{free\_page}：物理地址 → 计算 bit 索引 → 清 0
  \item 优点：简单，$O(1)$ free
  \item 缺点：alloc 是 $O(n)$，不支持连续多页分配
\end{itemize}

\section{Buddy 分配器}

Linux 风格的伙伴系统，支持 order 0..10（$2^0$ 到 $2^{10}$ 页）：

\begin{itemize}
  \item 每个 order 维护一个 free list
  \item \texttt{alloc}：找最小满足的 order → 如果没有则从更高 order 分裂
  \item \texttt{free}：归还后检查 buddy 是否也空闲 → 向上合并
\end{itemize}

\begin{thinkbox}
给定页帧号 $P$ 和分配的 order $k$，它的 buddy 页帧号是多少？
提示：XOR 运算——\texttt{buddy = P \^{} (1 << k)}。想想为什么 XOR 能找到 buddy。
\end{thinkbox}

\section{后端切换：\texttt{kconfig.h}}

\begin{lstlisting}[style=cstyle]
#define KCONFIG_PMM_BACKEND  0   // 0=bitmap, 1=buddy
\end{lstlisting}

\begin{lstlisting}[style=cstyle]
// kmain.c
#if KCONFIG_PMM_BACKEND == 0
    pmm_register_backend(pmm_bitmap_get_ops());
#elif KCONFIG_PMM_BACKEND == 1
    pmm_register_backend(pmm_buddy_get_ops());
#endif
\end{lstlisting}

\section{验证}

\begin{lstlisting}[style=shellstyle]
szy-kernel > pmm state
PMM: base=0x00100000, total=32768, free=32395, used=373
PMM: backend=bitmap
szy-kernel > pmm alloc 3
pmm alloc: 0x07FC1000
pmm alloc: 0x07FC2000
pmm alloc: 0x07FC3000
szy-kernel > pmm freeall
pmm freeall: freed 3 pages
\end{lstlisting}

\section{思考与练习}

\begin{exercise}
在 \texttt{kconfig.h} 中把 \texttt{KCONFIG\_PMM\_BACKEND} 改为 1（buddy），
重新编译运行。用 \texttt{pmm state} 对比 bitmap 和 buddy 的 free pages 数量。
两者为什么可能不同？（提示：buddy 的内部碎片和对齐需求。）
\end{exercise}

\begin{exercise}
运行 \texttt{pmm alloc 100} 然后 \texttt{pmm state}，观察 free 减少了多少。
再 \texttt{pmm freeall} 后 \texttt{pmm state}，free 是否完全恢复？
如果没有完全恢复，说明可能有内存泄漏——排查。
\end{exercise}

\begin{exercise}
在 buddy 分配器中，尝试实现 \texttt{alloc\_pages(order)} 接口：分配 $2^{\text{order}}$
个连续物理页。思考：bitmap 分配器能轻松支持这个需求吗？
\end{exercise}

\section{本章小结}

PMM 是内存管理的基石。可插拔后端模式让我们能在不改调用方的前提下
切换分配算法。下一章在 PMM 之上构建虚拟内存（分页）。
